Considerando o trade-off apresentado, 
foram estabelecidas hipóteses com o objetivo de verificar 
se a utilização de determinada abordagem apresenta vantagens 
quando observadas sob o mesmo cenário.
As hipóteses orientam a definição dos experimentos e a 
análise dos resultados obtidos.

\begin{itemize}
	\item \textbf{H1}: A quantização em inteiros de 8 bits, aplicada ao modelo implementado no fluxo de processamento em borda, causa redução de acurácia em comparação ao modelo quantizado em ponto flutuante de 32 bits, implementado no fluxo de processamento em nuvem.
    \item \textbf{H2}: O método de processamento local apresenta menor latência e maior valor de taxa de pacotes entregues, em comparação com o método de processamento remoto.
    \item \textbf{H3}: A abordagem de processamento em nuvem apresenta autonomia energética inferior durante a transmissão de mensagens quando comparada ao método de processamento local.
\end{itemize} 